\section{Capítulo 6 --- Dashboard, storytelling e decisão}

Uma rede de 24 clínicas acompanha acesso, qualidade e sustentabilidade. O painel mensal será usado pela diretoria para decidir onde abrir horários e onde investigar deterioração clínica.

\subsection{Questões objetivas}

\ItemHeader{1}{Objetiva}{hierarquia visual}{Situação profissional e suporte quantitativo}

\TextoBase

A página inicial do painel das 24 clínicas exibe 28 cartões com o mesmo tamanho, oito cores e quatro filtros abertos. Em dez segundos, a diretoria precisa localizar espera elevada sem aceitar uma redução obtida à custa de mais cancelamentos. A primeira camada deve priorizar a decisão; detalhes permanecem disponíveis para exploração. Na reunião, a decisão imediata é escolher duas unidades para abrir horários, mas a diretoria só autoriza a medida quando espera e cancelamento aparecem juntos, com período e meta visíveis. O protocolo também exige comparação com o mês anterior e indicação de qualidade dos dados. Assim, destaque visual não pode depender apenas de cor nem transformar um indicador isolado em recomendação; a leitura inicial precisa revelar prioridade, contraponto e próximo passo.

\FonteDidatica

\Comando Para orientar o redesenho da primeira camada, selecione a alternativa que preserva hierarquia, contexto e o indicador de qualidade.

\begin{enumerate}[label=\Alph*)]

\item ordenar todos os cartões alfabeticamente

\item usar uma cor diferente para cada unidade

\item priorizar uma Big Idea, destacar espera junto de cancelamento e relegar diagnóstico detalhado a uma camada de exploração

\item remover cancelamento para simplificar a mensagem

\item substituir números por um texto persuasivo

\end{enumerate}

\ItemHeader{2}{Objetiva}{média ponderada}{Situação profissional e suporte quantitativo}

\TextoBase

O memorando da diretoria apresenta no Quadro 1 tempos médios e volumes de três clínicas. O indicador único deve representar os 160 atendimentos, preservando o peso de cada volume.
\begin{DocumentBox}[Quadro 1 --- espera e volume]
\begin{tabularx}{\textwidth}{@{}Yrr@{}}\toprule
\textbf{Clínica} & \textbf{Espera média} & \textbf{Atendimentos}\\\midrule
A & 12 min & 100\\ B & 20 min & 50\\ C & 35 min & 10\\\bottomrule
\end{tabularx}
\end{DocumentBox}

\FonteDidatica

\Comando Calcule a espera média ponderada dos 160 atendimentos e selecione a alternativa correspondente.

\begin{enumerate}[label=\Alph*)]

\item 15,9 min, calculando a média ponderada pelo volume

\item 22,3 min, média simples das unidades

\item 67 min, soma das médias

\item 18,0 min, ponderando apenas as duas maiores

\item 20,0 min, usando a mediana sem considerar volume

\end{enumerate}

\ItemHeader{3}{Objetiva}{Big Idea}{Situação profissional e suporte quantitativo}

\TextoBase

Um piloto ampliou a triagem em parte das clínicas. A espera caiu apenas nas unidades participantes, enquanto o cancelamento permaneceu estável; nas demais, não houve melhora. Como alocará equipe no próximo mês, a diretoria precisa de uma mensagem central que preserve a condição observada e proponha monitoramento.

\FonteDidatica

\Comando Selecione a Big Idea que traduz a evidência do piloto em recomendação proporcional e verificável.

\begin{enumerate}[label=\Alph*)]

\item afirmar que triagem sempre reduz espera

\item esconder unidades sem melhora

\item recomendar expansão porque o gráfico ficou verde

\item formular mensagem condicionada: ampliar piloto de triagem, monitorando espera e cancelamento

\item trocar a hipótese por uma lista de indicadores

\end{enumerate}

\ItemHeader{4}{Objetiva}{taxa e incerteza}{Situação profissional e suporte quantitativo}

\TextoBase

Duas clínicas serão priorizadas para revisão de alta. A unidade A registrou 18 retornos em 200 altas; B, 12 em 80. Como os volumes diferem, a meta de reduzir retorno exige comparar taxas com seus denominadores, sem classificar pela contagem absoluta.

\FonteDidatica

\Comando Calcule as taxas de retorno de A e B e selecione a interpretação compatível com os tamanhos observados.

\begin{enumerate}[label=\Alph*)]

\item A=18\% e B=12\%, usando contagens como percentuais

\item A=9\% e B=15\%; B merece investigação, sem ignorar tamanhos amostrais

\item A=6\% e B=4\%, dividindo por 300

\item B é melhor porque teve menos retornos absolutos

\item as unidades são equivalentes porque a diferença é seis casos

\end{enumerate}

\ItemHeader{5}{Objetiva}{interação}{Situação profissional e suporte quantitativo}

\TextoBase

O filtro de período altera todos os gráficos, mas título, fonte e exportação não exibem o recorte ativo. Capturas do painel são levadas a reuniões e podem circular fora do sistema. A interação precisa preservar estado e permitir que outra pessoa reconstrua a leitura.

\FonteDidatica

\Comando Selecione o desenho de interação que torna o recorte reproduzível dentro e fora do painel.

\begin{enumerate}[label=\Alph*)]

\item ocultar filtros para evitar erro

\item usar animação como registro de estado

\item permitir qualquer combinação sem aviso

\item fixar sempre o último período escolhido

\item tornar estado visível no título, preservar filtros no artefato e oferecer restauração

\end{enumerate}

\ItemHeader{6}{Objetiva}{cauda}{Situação profissional e suporte quantitativo}

\TextoBase

Antes de uma mudança, o tempo mediano era 9 min, o p90 31 min e o p95 48 min. Depois, passaram a 8, 38 e 61 min, no mesmo tipo de atendimento. A direção precisa avaliar simultaneamente experiência central e cauda, evitando usar apenas a melhora de um minuto.

\FonteDidatica

\Comando Compare as três medidas antes e depois e selecione a conclusão sustentada pelo centro e pela cauda.

\begin{enumerate}[label=\Alph*)]

\item o serviço melhorou 11\% em todas as observações

\item o p95 melhorou 13 min

\item a mediana melhorou 1 min, mas a cauda piorou; não declarar melhora global

\item a média necessariamente caiu

\item percentis não podem coexistir em um dashboard

\end{enumerate}

\ItemHeader{7}{Objetiva}{acessibilidade}{Situação profissional e suporte quantitativo}

\TextoBase

O painel codifica meta cumprida em verde e falha em vermelho, sem texto, forma ou padrão redundante. Quando impresso em cinza, ambos os estados ficam visualmente semelhantes. O relatório será usado por públicos diversos e precisa manter o significado independentemente da percepção de cor.

\FonteDidatica

\Comando Selecione a revisão que mantém a distinção dos estados para impressão e diferentes percepções de cor.

\begin{enumerate}[label=\Alph*)]

\item aumentar saturação de vermelho e verde

\item combinar cor com texto, forma ou padrão e verificar contraste

\item adicionar mais tons intermediários

\item usar legenda distante sem rótulos

\item manter porque usuários aprenderão as cores

\end{enumerate}

\ItemHeader{8}{Objetiva}{variação}{Situação profissional e suporte quantitativo}

\TextoBase

A taxa de cancelamento em cinco semanas foi 8, 9, 8, 10 e 25\%. A média é 12\% e a mediana 9\%. A quinta semana coincide com falha de agendamento confirmada, mas o valor é legítimo; a equipe precisa comunicá-lo sem apagá-lo nem inventar tendência.

\FonteDidatica

\Comando Compare média, mediana e sequência semanal e selecione a ação exploratória proporcional ao pico.

\begin{enumerate}[label=\Alph*)]

\item declarar estabilidade porque a mediana é 9\%

\item remover 25\% sem registrar critério

\item usar somente média porque inclui tudo

\item destacar o pico e investigar sua causa; média isolada esconde a semana atípica

\item concluir tendência crescente com cinco valores

\end{enumerate}

\ItemHeader{9}{Objetiva}{storytelling}{Situação profissional e suporte quantitativo}

\TextoBase

Uma apresentação à diretoria começa pela ferramenta, percorre 16 gráficos e somente no final menciona a decisão sobre novos horários. Evidências, limites e recomendação aparecem separados, obrigando o público a reconstruir sozinho o argumento. O roteiro precisa servir à decisão, não demonstrar a interface.

\FonteDidatica

\Comando Selecione a reorganização narrativa que conecta contexto, evidência, limite e recomendação ao público decisor.

\begin{enumerate}[label=\Alph*)]

\item reordenar contexto, problema, evidência, limite e recomendação para o público decisor

\item começar por todas as transformações técnicas

\item eliminar limites para ganhar clareza

\item usar mais transições visuais

\item substituir evidência por depoimento

\end{enumerate}

\ItemHeader{10}{Objetiva}{comparação}{Situação profissional e suporte quantitativo}

\TextoBase

A meta de espera é 15 minutos. A clínica A registra espera 14 e cancelamento 22\%; B registra espera 17 e cancelamento 5\%, sob períodos comparáveis. Cada unidade atende melhor a um critério; a direção ainda não definiu qual risco tem prioridade.

\FonteDidatica

\Comando Avalie conjuntamente espera e cancelamento e selecione a conclusão que não agrega escalas nem inventa prioridade.

\begin{enumerate}[label=\Alph*)]

\item classificar A como melhor apenas pela espera

\item classificar B como melhor apenas por cancelamento

\item tirar média das duas métricas

\item ocultar a meta para evitar viés

\item mostrar conjuntamente espera e cancelamento; nenhuma unidade domina sem critério de prioridade

\end{enumerate}

\subsection{Questões discursivas}

\ItemHeader{11}{Discursiva}{Integrar evidências e justificar decisão}{Caso profissional do capítulo}

\TextoBase

A diretoria recebe uma tela única com espera, cancelamento, retorno e satisfação de 24 clínicas. Indicadores principais competem com filtros e diagnósticos, e capturas perdem o período selecionado. A equipe precisa separar decisão executiva de exploração sem romper a rastreabilidade.

\FonteDidatica

\Comando Em até 15 linhas, descreva duas camadas do dashboard. Declare Big Idea e público, organize hierarquia e interações e associe uma regra de decisão a espera e cancelamento. Explique como o recorte permanece auditável.

\ItemHeader{12}{Discursiva}{Integrar evidências e justificar decisão}{Caso profissional do capítulo}

\TextoBase

Três clínicas registraram espera média de 12, 20 e 35 minutos em 100, 50 e 10 atendimentos. Para o período consolidado, a mediana foi 9, o p90 38 e o p95 61 minutos. A diretoria quer representar volume, centro e cauda sem misturar suas funções.

\FonteDidatica

\Comando Calcule a média ponderada da rede, interprete os percentis e proponha uma combinação visual que mantenha visíveis volume e cauda. Explicite por que a média simples das clínicas seria inadequada.

\ItemHeader{13}{Discursiva}{Integrar evidências e justificar decisão}{Caso profissional do capítulo}

\TextoBase

Um piloto de triagem reduziu espera somente nas clínicas participantes e não alterou cancelamento. A apresentação atual começa pela ferramenta, mostra 16 gráficos e encerra sem decisão. O público precisa compreender evidência, alcance, incerteza e próximo passo em poucos minutos.

\FonteDidatica

\Comando Escreva um storyboard de cinco movimentos, do contexto à decisão. Em cada movimento, indique a função narrativa e inclua ao menos uma evidência, um limite e uma recomendação condicionada para ampliar ou interromper o piloto.

\ItemHeader{14}{Discursiva}{Integrar evidências e justificar decisão}{Caso profissional do capítulo}

\TextoBase

O painel usa vermelho e verde sem rótulos, possui fonte pequena e apresenta 28 cartões equivalentes. Em teste informal, usuários confundiram o período ativo e escolheram clínicas diferentes diante da mesma pergunta. A publicação será impressa e também navegada por tecnologias assistivas.

\FonteDidatica

\Comando Defina um teste de acessibilidade e um teste de interpretação, cada um com tarefa, participantes ou ferramenta, critério observável de aprovação e ação corretiva após falha.

\ItemHeader{15}{Discursiva}{Integrar evidências e justificar decisão}{Caso profissional do capítulo}

\TextoBase

Após uma mudança de agenda, a média do tempo caiu, mas o p95 aumentou 13 minutos. O painel declara ``melhora'' sem mostrar distribuição, canal, volume ou semanas individuais. A gestão considera expandir a mudança para todas as clínicas no mês seguinte.

\FonteDidatica

\Comando Em até 15 linhas, audite a conclusão: proponha hipóteses para centro e cauda, cortes de investigação, duas visualizações e uma regra mensurável para expandir, manter o piloto ou reverter.

\newpage

\subsection{Respostas explicadas das questões pares}

\paragraph{Questão 2 --- resposta A.} A soma ponderada é $12\times100+20\times50+35\times10=2.550$ minutos. Dividindo pelos $100+50+10=160$ atendimentos, obtém-se $2.550/160=15{,}9375$, ou 15,9 min. A média simples das três clínicas daria 22,3 min e super-representaria a unidade de apenas dez atendimentos.
\[\bar x_w=\frac{12(100)+20(50)+35(10)}{160}=15{,}9\text{ min}\]

\paragraph{Questão 4 --- resposta B.} Na unidade A, $18/200=0{,}09=9\%$; na B, $12/80=0{,}15=15\%$. Embora B tenha menos retornos absolutos, sua taxa é seis pontos percentuais maior. Ela merece investigação, preservando a informação de que seu denominador é menor. Isso corresponde à alternativa B.

\paragraph{Questão 6 --- resposta C.} A mediana caiu de 9 para 8 min, melhora de um minuto no centro. Entretanto, o p90 subiu de 31 para 38 min e o p95 de 48 para 61 min, pioras de 7 e 13 min na cauda. Portanto, não há base para declarar melhora global; a alternativa C integra as três medidas.

\paragraph{Questão 8 --- resposta D.} A média é $(8+9+8+10+25)/5=12\%$, enquanto a mediana é 9\%. A diferença decorre sobretudo do pico legítimo de 25\%, coincidente com a falha confirmada. Excluí-lo apagaria um evento operacional; usar apenas a média ocultaria a forma da série. A ação adequada é destacá-lo e investigar sua causa, como em D.

\paragraph{Questão 10 --- resposta E.} A atende à meta de espera por um minuto, mas tem 22\% de cancelamento; B excede a meta em dois minutos, porém tem apenas 5\% de cancelamento. Cada clínica vence em um critério. Sem prioridade ou função de decisão definida, não se deve fabricar um ranking: ambos os indicadores devem aparecer juntos. Logo, E.

\paragraph{Questão 12 --- critérios.} A resposta deve recalcular $2.550/160=15{,}9$ min, explicar por que a média simples é inadequada e comparar média e mediana da série de cancelamentos, identificando o pico de 25\%. O wireframe deve exibir volumes, período, meta, estado dos filtros e anotação do evento operacional.

\paragraph{Questão 14 --- critérios.} Espera-se roteiro que comece pela decisão, apresente evidências centrais e de cauda, explicite a falha da quinta semana e compare espera e cancelamento sem prioridade inventada. A recomendação deve ser proporcional e acompanhada de critério de monitoramento e fonte reproduzível.
